\section{Reserved and Compatibility Rules}
This section defines the default compatibility contract for reserved fields, optional
features, and future architectural growth. A later section may narrow a rule for a
specific structure, but it must say so explicitly.

\subsection{Terminology}
\begingroup\footnotesize
\begin{longtable}{@{}p{1.45in}p{4.05in}@{}}
\toprule
\textbf{Term} & \textbf{Meaning}\\
\midrule
\endhead
\texttt{reserved} & Architecturally held for a future definition. Software must not depend on its current value or behavior.\\
\texttt{must be zero} & Software must write zero. A nonzero value is invalid unless a later architectural revision assigns the field.\\
\texttt{read-as-zero} & Hardware returns zero for the field on this implementation. Software must still treat the field as reserved unless it is explicitly software-defined.\\
\texttt{ignored} & Hardware accepts the field but does not use it for the current operation. Ignored fields are not automatically available for software-defined state.\\
\texttt{preserved} & Software must retain the previous value when rewriting the containing object.\\
\texttt{software-defined} & Hardware does not interpret the field and reserves it for software use. Future architectural extensions do not consume this field without redefining the containing structure.\\
\texttt{implementation-defined} & The implementation chooses the value or behavior and exposes the relevant contract through CPUID, fixed documentation, or platform firmware.\\
\texttt{illegal} & Use of a disallowed instruction opcode encoding, effective-address form, selector, or operation raises the exception named by the defining rule. Prefix bytes use the prefix model's unassigned-value rule.\\
\bottomrule
\end{longtable}
\manualtablecaption{Compatibility Terms}
\endgroup

\subsection{Reserved Fields}
Architected registers and architectural memory structures use the following defaults:

\begin{itemize}
\item Reserved bits in architected registers read as zero and must be written as zero.
\item Reserved bits in control registers or selector fields cause \texttt{INVALID\_CONTROL\_STATE} when a nonzero or reserved value is written.
\item Reserved bits in page-table entries cause \texttt{PAGE\_FAULT} when they are consumed by page-table translation.
\item Reserved bits in interrupt-vector-table entries and supervisor-entry frames must be zero. A reserved CPU vector has no predefined frame type; attempted delivery follows the double-fault rule described by the exception-processing model.
\item Software-defined fields are not reserved fields. Software may store values in them, and hardware ignores them except where the containing structure explicitly says otherwise.
\end{itemize}

\subsection{Reserved Encodings}
Reserved instruction opcodes, reserved extension opcode encodings,
reserved effective-address forms, and unsupported optional instruction-group encodings raise
the exception named in the compatibility policy table. Bedrock keeps the decode
contract simple: if the implementation does not support an optional instruction
group, encodings from that group are not distinguished from other illegal
instruction encodings by a separate extension-availability vector.

Prefix bytes are not instruction opcodes. Unassigned prefix values follow the
prefix policy values in the compatibility policy table and remain reserved for
future prefix definitions. Prefix bytes are consumed only as prefix metadata;
instruction decode then proceeds from the following opcode, operands, and
addressing forms.

Noncanonical encodings are invalid unless the instruction description explicitly defines an
alias or priority rule. Assemblers must emit canonical encodings by default. Disassemblers
must prefer canonical spellings when more than one spelling names the same operation.

\subsection{CPUID Compatibility}
\texttt{CPUID} is the compatibility boundary for optional architectural behavior.
Software must not use an optional instruction group, optional register state, or optional
state image unless the corresponding CPUID bit or leaf reports support.

Unknown CPUID selectors, reserved CPUID result bits, privilege, serialization, and
runtime mutability are governed by the compatibility policy table unless a CPUID
leaf explicitly gives a narrower rule.

@COMPATIBILITY_POLICY_TABLE@
